\documentclass[10pt]{article}
\usepackage[margin=20mm]{geometry}
\usepackage{amsmath,amssymb,fontspec,unicode-math,xurl,hyperref}
\IfFontExistsTF{DejaVu Serif}{\setmainfont{DejaVu Serif}}{\setmainfont{Latin Modern Roman}}
\setmathfont{latinmodern-math.otf}
\setlength{\parindent}{0pt}\setlength{\parskip}{6pt}\sloppy
\begin{document}
Position: LLM Social Simulations Are a Promising Research Method\par
Claim-focused scientific review ICML 2025 Human decision pending
Paper I25-40125 | Report 0023 | AI review; human decision pending\par
---\par
\section*{1. Identity, reviewed version and outcome}
Jacy R. Anthis and colleagues. ICML 2025 position paper, PMLR 267:81005-81034. Reviewed the publisher-linked 30-page PDF retrieved on 24 October 2026 UTC; all target locators below use printed internal pages. OpenReview ID: cRBg1dtj7o. Source SHA-256: `75135835b88ceb84aab6ec62cbcf56a06e8a77cea07650a334f67ec54c90dd31`. [S1]\par
Supported outcome: Minor. A local statistical interpretation problem affects prominent evidence: correlations and a likelihood-based pseudo-R-squared are described as proportions of variation predicted. Correcting the metric names and interpretations leaves the central claim of a promising, conditional research direction intact. No Major issue is supported. This review does not establish the validity of simulations for any untested population or intervention.\par
\section*{2. Inspected scope and strongest alternative reading}
Read the main argument and Appendices A-C, and checked the three prominent empirical examples against dated primary sources. Table A1 has 53 rows across pp. 22-25 (15, 16, 14 and 8); Table A2 has 10 selected rows on p. 26. The tables contain no performance-value column. Their row counts are consistent with their framing, but that does not establish exhaustive literature coverage or independence of studies.\par
The position is prospective and explicitly conditional. Its five challenges, validation recommendations, limitations of the narrative literature search and emphasis on individual-response simulation should constrain interpretation. The paper does not claim that all current simulations can replace human experiments. Its discussion of imperfect transfer and preregistered validation is material counterevidence to a broader accusation of ignoring generalization. This review therefore distinguishes a metric error from disagreement with the authors’ research priorities.\par
\section*{3. Problem 1 - Correlation and likelihood fit are not explained variation}
Source claims and precise locations. The Introduction, p. 1 right column, describes the Hewitt result as 91\% of variation predicted; §4.5.2, p. 8 right column, similarly renders adjusted correlations as percentages of variation. Appendix A.2.2, beginning on p. 28 and continuing on p. 29, describes the Centaur performance range in terms of variation in human data. The invalid step is assigning a variance-explanation interpretation to metrics that measure different quantities.\par
Primary-source check for Hewitt. The dated 8 August 2024 working paper defines Pearson correlation and its disattenuated counterpart on PDF p. 5; Figure 2, p. 7, reports the comparisons. GPT-4 has raw r = 0.85 and adjusted r = 0.91. The adjusted comparison values are -0.09, 0.25, 0.82, 0.84 and 0.91, matching the target’s percentages. These are correlations, not calibrated prediction R-squared values. The target itself correctly labels r = 0.91 in Appendix A.2.4, p. 29. [S1, S2]\par
Independent derivation. Let Y and P be centered random variables with unit variance, finite second moments and correlation rho. For a linear calibrated predictor aP, the mean squared error is the following quadratic. Its derivative vanishes at a = rho, and its second derivative is 2, so this is the unique minimum.\par
\[\mathbb{E}(Y-aP)^2=1-2a\rho+a^2,\qquad a^*=\rho,\qquad R^2_{\rm calibrated}=1-(1-\rho^2)=\rho^2.\] Thus even in this favorable, fully specified calibration example, rho = 0.91 gives R-squared = 0.8281, not 0.91. More generally, a forecast’s out-of-sample R-squared also depends on its scale, offset, comparison baseline and evaluation set. This calculation is an illustration of the invalid inference, not a reanalysis of the source data. In particular, 82.81\% must not be substituted as an empirical predictive-variance result: the source’s 0.91 is disattenuated and is not the correlation of a supplied calibrated prediction vector.\par
Calibration is a substantive distinction here: Hewitt’s PDF p. 8 reports systematic overestimation of effect magnitudes and evaluates a rescaling correction. This does not invalidate the reported association; it shows why association and numerical prediction accuracy should be separately identified. [S2]\par
Primary-source check for Centaur. Binz et al., arXiv:2410.20268v1, PDF p. 12 Eq. (1), defines a normalized log-likelihood metric relative to uniform guessing. PDF p. 4 names it pseudo-R-squared and describes a participant-level training/test split. The source’s 160-experiment dataset should not be interpreted as 160 unseen experiments. [S3]\par
\[R^2_{\rm pseudo}=1-\frac{\log p_{\rm model}(D)}{\log p_{\rm guess}(D)}.\] Complete metric-distinction example. Take independent Bernoulli responses with success probability p = 0.9. Let the model assign probability 0.9 to success on every trial and the guessing baseline assign 0.5. Both are valid probability distributions, all likelihoods are nonzero, and expected log likelihoods are finite. On a per-observation expected-log-likelihood basis:\par
\[H(p)=-p\log p-(1-p)\log(1-p),\quad R^2_{\rm pseudo}=1-\frac{H(0.9)}{\log 2}=0.5310044064\ldots\] The same model’s conditional-mean prediction is the constant 0.9. Its squared-error risk is 0.9(0.1) = 0.09, exactly Var(Y). Therefore its variance-based predictive R-squared relative to the mean baseline is zero. The positive pseudo-R-squared arises from improving over uniform probabilities; it does not mean 53.1\% of response variance is explained. This is a population calculation illustrating metric non-equivalence, not an estimate from Centaur’s experiments.\par
Counterevidence and severity. A charitable reading treats ’variation’ as informal language for statistical correspondence. Section 4.5.2 even explicitly mentions correlation, and the appendix supplies the correct r notation. Those qualifications make a central-invalidity claim unwarranted, but they do not make a numerical percentage-of-variation interpretation correct. The error concerns how supporting evidence is communicated; the high reported association and the motivation to investigate simulation remain. Judgment: Minor.\par
Repair and verification. Replace the Hewitt wording with a disattenuated correlation of 0.91, optionally accompanied by the raw correlation of 0.85, and label all comparisons as correlations. Describe Centaur using likelihood-based predictive pseudo-R-squared relative to uniform guessing. Do not translate either result into a variance percentage without a separately defined and evaluated variance-based metric. These wording repairs follow directly from the primary metric definitions and require no new experiment. Author adoption has not been requested.\par
\section*{4. Checks that did not support additional problems}
Claim / location | Check and bounded conclusion\par
Park’s 85\% headline; p. 1 and Appendix A.2.3, p. 29 | Historical v1, pp. 4 and 6, defines retest-normalized accuracy and reports 0.85; raw accuracy is 68.85\%, with human retest accuracy 81.25\%. The target preserves normalization. No Issue Found on this point. [S4]\par
Centaur’s 160 experiments; p. 1 | Historical v1, p. 12, confirms dataset size. Comparative performance and generalization analyses were not reproduced. [S3]\par
Literature tables and search scope; pp. 22-26, 29-30 | Row counts checked. Narrative selection and subjective challenge coding are disclosed. No unsupported claim of systematic-review completeness is inferred. [S1]\par
Validation and future feasibility; §4.5.2, p. 8 | Prospective recommendations and hypotheses, not proved universal guarantees. No counterexample to a universal theorem is claimed. [S1]\par
The Park record has later revisions, including a renamed June 2026 version. The historical 15 November 2024 v1 is the appropriate comparison for the target’s 2024 citation; later values are not used to manufacture a retrospective discrepancy. The target’s own arXiv record, 2504.02234, lists v1 (3 April 2025) and v2 (5 June 2025). This review concerns the PMLR PDF; no byte-equivalence or exhaustive erratum-history claim is made. [S4, S5]\par
\section*{5. Coverage limits and decision status}
No simulation model, author code or human-subject experiment was run. The review did not re-estimate Hewitt correlations, refit Centaur, reproduce Park’s individual-level ratios, reconstruct the entire literature search, audit every cited paper, or validate all OOD examples. The attached coverage ledger separates these unperformed tasks from the completed source and mathematical checks. No separate author code or supplement was identified on the publisher record beyond the embedded appendices. OpenReview discussion history was not verified.\par
Historical primary sources, table counts and locators were independently checked; the initial p. 7 locator was corrected to p. 8 against the exact PDF. Human scientific decision: pending. Author adoption: not requested.\par
\section*{References and evidence locations}
[S1] Anthis et al. (2025), reviewed PMLR267 position paper, pp. 81005-81034. Download record and exact PDF retained.
https://proceedings.mlr.press/v267/anthis25a.html\par
[S2] Hewitt et al., Predicting Results of Social Science Experiments Using Large Language Models, dated 8 August 2024, 18-page primary working paper. Retrieved from a mirror after the OSF route was unavailable; SHA-256 `7258b046dfe22a6d7eb1e1f76cc0a2f4fd7473f571cc259ddfa386b01a738aa7`.
https://samim.io/dl/Predicting\%20results\%20of\%20social\%20science\%20experiments\%20using\%20large\%20language\%20models.pdf\par
[S3] Binz et al., A Foundation Model of Human Cognition, arXiv:2410.20268v1, 141 pages. Selected pp. 4 and 12 inspected; PDF SHA-256 `b5a41f69b522386ef4c2798be6361d9b3ccdcaaad892529e5ea5b98bcb3951e8`.
https://arxiv.org/pdf/2410.20268v1\par
[S4] Park et al., Generative Agent Simulations of 1,000 People, arXiv:2411.10109v1, 15 November 2024, 65 pages. Selected pp. 4 and 6 inspected; PDF SHA-256 `719ece4577b62e2b8cdafb35101455051c50bda1b4c8e9da5d0ed339e96cc348`.
https://arxiv.org/abs/2411.10109v1\par
[S5] Target arXiv version metadata, accessed 24 October 2026 UTC. Metadata inspection only; PMLR PDF remains the reviewed version.
https://arxiv.org/abs/2504.02234\par
\section*{Supporting checks}
{\ttfamily\footnotesize python sources/0023/independent\_checks.py\\
calibrated R2 = 8281/10000; pseudo-R2 = 0.5310044064107189\\
Bernoulli variance-based R2 = 0}\par
Source p. 8 was visually inspected.\par
\end{document}
